% ===================================================================
%  TABLICA Nr 15 — Targ. — Artykuły spożywcze.
%  Meme regime que les precedents. Tableau a UNE SEULE COLONNE.
%
%  UN RENVOI DE CE TABLEAU EST CORRIGE SUR LA PAGE DE LECTURE, DANS
%  LA COLONNE FRANCAISE SEULE : elle imprime (32) au 13 la ou l'ido
%  lit (82), les petits pains. La correction vit dans
%  gravuri/korekti.json et ne touche aucune source. Cette colonne
%  ecrit le numero de l'ido.
%
%  LE TABLEAU EST UN INVENTAIRE, et c'est la sa difficulte : cent
%  cinq renvois, presque tous des denrees, et beaucoup se ressemblent
%  assez pour tomber sur le meme mot polonais si on n'y prend pas
%  garde. « prunes » (53) et « prunes seches » (65) sont deux
%  renvois ; « korbi » (70) et « provizo-reti » (73) aussi ; et le
%  tableau distingue « bucho-servisto » (75), le garcon qui pousse la
%  brouette, de « bucho-garsono » (93), celui qui decoupe.
%
%  ET LES TROIS VIANDES DOIVENT RESTER TROIS : « muton-karno » (90),
%  « bov-karno » (94) et « bovyun-karno » (92) — baranina, wolowina,
%  cielecina. Le polonais a les trois mots ; une colonne qui dirait
%  « mieso wolowe » deux fois perdrait un renvoi de la planche.
% ===================================================================

%%K t15-apar-1 apar
\VUsaut{4.0mm}
\VUtitre{15.0pt}{60}{TABLICA Nr 15}
\VUsaut{3.0mm}

%%K t15-tit-1 sub
\VUpk{11.4pt}{40}{Targ. --- Artykuły spożywcze.}
\VUsaut{3.0mm}

%%K t15-01-1 p
1. --- Większość kupców, którzy dostarczają nam towarów potrzebnych do naszego żywienia,
zbiera się pod \VUgras{halą} \textsuperscript{(1)} \VUgras{targu krytego} albo na
sąsiednich ulicach.

%%K t15-02-1 p
2. --- \VUgras{Wędliniarz} \textsuperscript{(2)}, który sprzedaje \VUgras{szynkę}
\textsuperscript{(3)}, \VUgras{kaszankę} \textsuperscript{(4)}, \VUgras{kiełbaskę}
\textsuperscript{(5)}, \VUgras{kiszkę} \textsuperscript{(6)} i \VUgras{słoninę}
\textsuperscript{(7)}, stoi blisko \VUgras{sprzedawczyni masła i sera}
\textsuperscript{(8)}, której wystawa jest zaopatrzona w \VUgras{sery holenderskie}
\textsuperscript{(9)} z czerwoną skórką, w \VUgras{sery camembert}
\textsuperscript{(10)}, w \VUgras{krążki gruyère} \textsuperscript{(11)} i w świeże
\VUgras{bloki masła} \textsuperscript{(12)}.

%%K t15-03-1 p
3. --- Przed nią \VUgras{sprzedawczyni jarzyn} \textsuperscript{(13)} proponuje swoim
klientkom piękne \VUgras{pomidory} \textsuperscript{(14)}, \VUgras{kalafiory}
\textsuperscript{(15)}, \VUgras{sałatę} \textsuperscript{(16)}, \VUgras{kapustę
głowiastą} \textsuperscript{(17)}, \VUgras{kabaczki} \textsuperscript{(19)},
\VUgras{melony} \textsuperscript{(18)}, a nawet \VUgras{grzyby} \textsuperscript{(20)}.
Ale \VUgras{roznosiciel piwa}, który zatrzymał swój \VUgras{wóz dostawczy}
\textsuperscript{(21)} obładowany \VUgras{beczkami piwa} \textsuperscript{(22)} właśnie
przed jej wystawą, przeszkadza jej klientkom się zbliżyć. Ogrodniczka przynosi jej
wielki kosz \VUgras{karczochów} \textsuperscript{(23)}.

%%K t15-tit-2 sub
\VUsaut{4.0mm}
\VUpk{10.2pt}{40}{Sprzedaż i kupno.}
\VUsaut{3.0mm}

%%K t15-04-1 p
4. --- Na targu ożywienie jest wielkie, bo nadeszła godzina \VUgras{sprzedaży
licytacyjnej} \textsuperscript{(24)}. \VUgras{Gospodynie domowe} \textsuperscript{(26)},
które przychodzą tam po swoje zakupy, zatrzymują się przechodząc przed straganem
\VUgras{sprzedawczyni flaków} \textsuperscript{(25)}, żeby kupić \VUgras{flaki} albo
\VUgras{głowę cielęcą.}

%%K t15-05-1 p
5. --- Tłusty \VUgras{kucharz} \textsuperscript{(27)} targuje bardzo pięknego
\VUgras{tuńczyka} u \VUgras{sprzedawczyni ryb} \textsuperscript{(28)}, która według
swojego uznania podnosi swoje ceny sprzedaży. Dostała wielką ilość \VUgras{węgorzy,
soli, pstrągów} i \VUgras{karpi.} Jej \VUgras{dorsz} i \VUgras{małże} są bardzo świeże.

%%K t15-tit-3 sub
\VUsaut{4.0mm}
\VUpk{10.2pt}{40}{Tanie i drogie towary.}
\VUsaut{3.0mm}

%%K t15-06-1 p
6. --- \VUgras{Sprzedawczyni drobiu} \textsuperscript{(29)}, umieszczona blisko niej,
jest bardzo sumienna. Kiedy sprzedaje \VUgras{indyka, gęś, kaczkę} albo zwykłe
\VUgras{kurczę,} żąda tylko słusznej ceny.

%%K t15-07-1 p
7. --- Na rogu placu, na pierwszym piętrze, od niedawna urządził się \VUgras{sprzedawca
płótna} \textsuperscript{(30)}, u którego kupujące są zawsze pewne, że znajdą bardzo
piękny wybór \VUgras{obrusów,} \VUgras{serwetek, fartuchów} i \VUgras{ręczników.} Na tym
samym piętrze \VUgras{nożownik} \textsuperscript{(31)} urządził swoją pracownię. Ostrzy
\VUgras{noże} sprzedawczyń w hali i wyrabia dla ogrodników \VUgras{sierpy} z
najtwardszej \VUgras{stali.}

%%K t15-tit-4 sub
\VUsaut{4.0mm}
\VUpk{10.2pt}{40}{Sklep korzenny.}
\VUsaut{3.0mm}

%%K t15-08-1 p
8. --- Pod jego pracownią od niedawna otwarto wielki magazyn artykułów spożywczych. Już
nie jest to zwykły i nieważny dawny \VUgras{sklep korzenny} \textsuperscript{(32)},
który sprzedawał tylko towary powszechnego użytku, \VUgras{herbatę}
\textsuperscript{(33)}, \VUgras{czekoladę} \textsuperscript{(34)}, \VUgras{cukier w
głowach} \textsuperscript{(37)} i \VUgras{mydło} \textsuperscript{(38)}. Tam sprzedaje
się każdy towar, \VUgras{chrupiące ciastka}, \VUgras{konserwy} \textsuperscript{(35)},
\VUgras{likiery} \textsuperscript{(36)}, \VUgras{rum,} \VUgras{koniak, kirsz,
anyżówkę} i \VUgras{curaçao.} Tam znajduje się zwierzynę: \VUgras{zająca}
\textsuperscript{(39)}, \VUgras{drozdy} \textsuperscript{(40)}, \VUgras{kuropatwy}
\textsuperscript{(41)}, \VUgras{przepiórki} \textsuperscript{(42)}, \VUgras{dziką
kaczkę} \textsuperscript{(43)} albo \VUgras{bażanta} \textsuperscript{(44)}, równie
łatwo jak owoce egzotyczne takie jak \VUgras{banany} \textsuperscript{(46)} i
\VUgras{ananasy.}

%%K t15-09-1 p
9. --- Bardzo uprzejmy \VUgras{czeladnik} \textsuperscript{(45)} korzennika jest gotów
dać to, czego się pragnie: \VUgras{konfiturę} \textsuperscript{(47)} porzeczkową albo
pigwową, \VUgras{smalec} \textsuperscript{(48)}, \VUgras{ogórki konserwowe}
\textsuperscript{(49)} albo solone \VUgras{sardynki} \textsuperscript{(50)}.

%%K t15-09-2 p
To jest bardzo dogodne dla \VUgras{klientek} \textsuperscript{(51)}, które znajdują obok
towarów najpowszechniejszego użytku najrzadsze nowalijki i owoce najsmakowitsze:
soczyste \VUgras{gruszki} \textsuperscript{(52)}, \VUgras{śliwki} \textsuperscript{(53)},
\VUgras{winogrona} \textsuperscript{(54)}, \VUgras{brzoskwinie} \textsuperscript{(55)},
\VUgras{migdały} \textsuperscript{(56)} i \VUgras{orzechy} \textsuperscript{(57)}.

%%K t15-tit-5 sub
\VUsaut{4.0mm}
\VUpk{10.2pt}{40}{Klientela.}
\VUsaut{3.0mm}

%%K t15-10-1 p
10. --- \VUgras{Ryż} \textsuperscript{(58)} i \VUgras{soczewica} \textsuperscript{(59)}
są obok skrzyni wędzonych \VUgras{śledzi} \textsuperscript{(60)}; \VUgras{ziemniak}
\textsuperscript{(61)} i \VUgras{fasola} \textsuperscript{(63)} sąsiadują z
\VUgras{jabłkami} \textsuperscript{(62)}, \VUgras{figami} \textsuperscript{(64)},
\VUgras{suszonymi śliwkami} \textsuperscript{(65)} i \VUgras{orzechami laskowymi}
\textsuperscript{(66)}. Znajduje się w tym magazynie świeże \VUgras{jajka}
\textsuperscript{(67)} i \VUgras{oliwki} \textsuperscript{(68)}; dlatego \VUgras{kosze}
\textsuperscript{(70)} i \VUgras{siatki na zakupy} \textsuperscript{(73)} są napełniane
z ogromną szybkością.

%%K t15-11-1 p
11. --- \VUgras{Kawa} \textsuperscript{(72)} tej wybornej firmy zyskała wśród
\VUgras{służących} \textsuperscript{(69)} i \VUgras{kucharek} słuszną sławę, bo stary
\VUgras{korzennik} \textsuperscript{(71)} sam ją najstaranniej praży.

%%K t15-tit-6 sub
\VUsaut{4.0mm}
\VUpk{10.2pt}{40}{Widowiska.}
\VUsaut{3.0mm}

%%K t15-12-1 p
12. --- Nie wszyscy idą na targ, żeby się zaopatrzyć. W malowniczym ruchu na uliczce,
która prowadzi do hali, widzi się osoby najrozmaitsze. Obok uprzejmego \VUgras{pomocnika
rzeźnika} \textsuperscript{(75)}, który pcha przed sobą \VUgras{taczkę}
\textsuperscript{(74)}, oto dwaj żołnierze na przepustce, całkiem dumni, że pokazują
swój lśniący mundur: \VUgras{huzar} \textsuperscript{(76)} z niebieskim
\VUgras{dolmanem} i \VUgras{czakiem z pióropuszem}, i \VUgras{kapral zuawów,} z
bufiastymi spodniami i głową nakrytą \VUgras{fezem.}

%%K t15-13-1 p
13. --- \VUgras{Piekarz} \textsuperscript{(80)}, który stoi na progu \VUgras{piekarni}
\textsuperscript{(78)}, właśnie ugniótł ciasto swojego \VUgras{chleba}
\textsuperscript{(79)} i wyjął z pieca \VUgras{rogaliki} \textsuperscript{(81)},
\VUgras{bułki} \textsuperscript{(82)} i \VUgras{okrągłe chleby} \textsuperscript{(83)},
które są w witrynie.

%%K t15-14-1 p
14. --- Nad piekarnią mieszka zdolna \VUgras{bieliźniarka} \textsuperscript{(84)}, która
zatrudnia kilka \VUgras{hafciarek} \textsuperscript{(85)}. Obok niej mieszka
\VUgras{praczka-prasowaczka} \textsuperscript{(86)}, która krochmali \VUgras{bieliznę}
\textsuperscript{(88)} po upraniu jej i wysuszeniu, i która wygładza ją
\VUgras{żelazkiem} \textsuperscript{(87)}.

%%K t15-tit-7 sub
\VUsaut{4.0mm}
\VUpk{10.2pt}{40}{Mięso, ryby i jarzyny.}
\VUsaut{3.0mm}

%%K t15-15-1 p
15. --- W \VUgras{jatce} \textsuperscript{(89)}, położonej na parterze, sprzedaje się
wszelkiego rodzaju \VUgras{mięsa, baraninę} \textsuperscript{(90)}, \VUgras{wołowinę}
\textsuperscript{(94)} albo \VUgras{cielęcinę} \textsuperscript{(92)} jako \VUgras{udźce
baranie, befsztyki, pieczenie} i \VUgras{kotlety.} \VUgras{Czeladnik rzeźnicki}
\textsuperscript{(93)} rozbiera wszystkie te mięsa i osobno kładzie \VUgras{kości
szpikowe} \textsuperscript{(96)}. Przy drzwiach jatki jest \VUgras{sprzedawczyni ostryg}
\textsuperscript{(97)}, która sprzedaje \VUgras{ostrygi} \textsuperscript{(98)}. Otwiera
je, jeśli się tego pragnie.

%%K t15-16-1 p
16. --- Wszyscy ci kupcy płacą patent albo opłatę za miejsce; dlatego \VUgras{policjant}
\textsuperscript{(100)} bez litości ściga biedne \VUgras{przekupki jarzyn}
\textsuperscript{(99)}, które robią im konkurencję, roznosząc na swoich wózkach
\VUgras{marchew, rzodkiewki, rzepę, pory} albo \VUgras{pęczki szparagów, świeży groszek,
szpinak, szczaw, rzeżuchę} i \VUgras{pietruszkę.}

%%K t15-tit-8 sub
\VUsaut{4.0mm}
\VUpk{10.2pt}{40}{Uwaga na policjanta!}
\VUsaut{3.0mm}

%%K t15-17-1 p
17. --- Chłopiec, który sprzedaje \VUgras{czosnek} \textsuperscript{(101)} i
\VUgras{cebulę,} bardzo dobrze wie, że i jemu nie wolno sprzedawać swoich towarów blisko
targu. Ucieka, biegnąc z ryzykiem przewrócenia kosza \VUgras{krabów,} \VUgras{raków} i
\VUgras{krewetek} \VUgras{sprzedawcy} \textsuperscript{(102)} \VUgras{langust} i
\VUgras{homarów.}

%%K t15-18-1 p
18. --- Wszyscy się krzątają i hałasują; ale to nie przeszkadza słyszeć wyraźnie głosu
wędrownego \VUgras{szklarza} \textsuperscript{(103)} ani okrzyku \VUgras{węglarza}
\textsuperscript{(104)}, który właśnie zostawił na kilka chwil swój wózek, żeby
dostarczyć \VUgras{worek węgla} \textsuperscript{(105)}.
\VUsaut{6.0mm}
\VUfilet{22mm}
